% GENERATED FILE. Edit canonical lessons, then run npm run generate:books.
% canonical-source-hash: fnv1a64:448a3cfb60544544
% canonical-lessons: JA-C386-ryoken, JA-C386-zeikan, JA-C386-nyuukoku, JA-C386-shukkoku, JA-C386-toridasu

\chapter{Passport and Customs}
\label{ch:ja-passport-and-customs}
\hlchaptermodality{386}

\begin{chapteropening}
\textbf{By the end of this chapter:} \emph{I can get through passport control and customs.}

Complete the last lesson of chapter 386: I can get through passport control and customs.
\end{chapteropening}

\section[\texorpdfstring{\ja{りょけん} (ryoken)}{りょけん (ryoken)}]{\ja{りょけん} (ryoken) — a passport}

\label{lesson:JA-C386-ryoken}



\begin{quote}
Before the new one: say the Japanese for a time difference, then the Japanese for to aim for.
\end{quote}

\subsection*{You'll want to know: \ja{りょけん}}
\textbf{\ja{りょけん}} — \emph{ryoken} — \textquotedblleft{}a passport\textquotedblright{}.

\begin{quote}
\textbf{\ja{りょけんを} \ja{なくさないで} \ja{ください。}} — \emph{ryoken o nakusanaide kudasai} — \textquotedblleft{}Please do not lose your passport.\textquotedblright{}
\end{quote}

\begin{cousinweb}
\textbf{\ja{りょけん}} (\emph{ryoken}) is Sino-Japanese (kango): two Chinese characters read with their on'yomi, readings borrowed from Chinese. They are \emph{ryo}, \textquotedblleft{}trip\textquotedblright{} (Middle Chinese \emph{ljoX}) and \emph{ken}, \textquotedblleft{}ticket\textquotedblright{} (\emph{khjwonH}) (KANJIDIC2; Middle Chinese in Baxter's transcription).
\end{cousinweb}

\begin{grammarlens}[title={What you've built}]
One more word for getting around.
\end{grammarlens}

\subsection*{Your turn}
\begin{itemize}
\raggedright
  \item \emph{Say it:} \emph{ryoken}
  \item \emph{Say it:} \emph{ryoken}, once more
  \item \emph{Say it:} \emph{ryoken o nakusanaide kudasai} — Please do not lose your passport.
  \item \emph{Recall:} say the Japanese for a time difference, then the Japanese for to aim for, then say \emph{ryoken} again
\end{itemize}

\begin{tcolorbox}[breakable,colback=teal!4,colframe=teal!35!black,title={Before you move on}]
What does \ja{りょけん} mean? (\textquotedblleft{}A passport\textquotedblright{}.) Say it once more.
\end{tcolorbox}
\section[\texorpdfstring{\ja{ぜいかん} (zeikan)}{ぜいかん (zeikan)}]{\ja{ぜいかん} (zeikan) — customs}

\label{lesson:JA-C386-zeikan}



\begin{quote}
Before the new one: say the Japanese for to aim for, then the Japanese for a passport.
\end{quote}

\subsection*{You'll want to know: \ja{ぜいかん}}
\textbf{\ja{ぜいかん}} — \emph{zeikan} — \textquotedblleft{}customs\textquotedblright{}.

\begin{quote}
\textbf{\ja{くうこうで} \ja{ぜいかんを} \ja{とおりました。}} — \emph{kūkō de zeikan o tōrimashita} — \textquotedblleft{}I went through customs at the airport.\textquotedblright{}
\end{quote}

\begin{cousinweb}
\textbf{\ja{ぜいかん}} (\emph{zeikan}) is Sino-Japanese (kango): two Chinese characters read with their on'yomi, readings borrowed from Chinese. They are \emph{zei}, \textquotedblleft{}tax\textquotedblright{} (Middle Chinese \emph{sywejH}) and \emph{kan}, \textquotedblleft{}barrier\textquotedblright{} (\emph{kwaen}) (KANJIDIC2; Middle Chinese in Baxter's transcription).
\end{cousinweb}

\begin{grammarlens}[title={What you've built}]
One more word for getting around.
\end{grammarlens}

\subsection*{Your turn}
\begin{itemize}
\raggedright
  \item \emph{Say it:} \emph{zeikan}
  \item \emph{Say it:} \emph{zeikan}, once more
  \item \emph{Say it:} \emph{kūkō de zeikan o tōrimashita} — I went through customs at the airport.
  \item \emph{Recall:} say the Japanese for to aim for, then the Japanese for a passport, then say \emph{zeikan} again
\end{itemize}

\begin{tcolorbox}[breakable,colback=teal!4,colframe=teal!35!black,title={Before you move on}]
What does \ja{ぜいかん} mean? (\textquotedblleft{}Customs\textquotedblright{}.) Say it once more.
\end{tcolorbox}
\section[\texorpdfstring{\ja{にゅうこく} (nyūkoku)}{にゅうこく (nyūkoku)}]{\ja{にゅうこく} (nyūkoku) — entry into a country}

\label{lesson:JA-C386-nyuukoku}



\begin{quote}
Before the new one: say the Japanese for a passport, then the Japanese for customs.
\end{quote}

\subsection*{You'll want to know: \ja{にゅうこく}}
\textbf{\ja{にゅうこく}} — \emph{nyūkoku} — \textquotedblleft{}entry into a country\textquotedblright{}.

\begin{quote}
\textbf{\ja{にゅうこくの} \ja{とき、りょけんが} \ja{ひつようです。}} — \emph{nyūkoku no toki, ryoken ga hitsuyō desu} — \textquotedblleft{}You need your passport when you enter the country.\textquotedblright{}
\end{quote}

\begin{cousinweb}
\textbf{\ja{にゅうこく}} (\emph{nyūkoku}) is Sino-Japanese (kango): two Chinese characters read with their on'yomi, readings borrowed from Chinese. They are \emph{nyū}, \textquotedblleft{}enter\textquotedblright{} (Middle Chinese \emph{nyip}) and \emph{koku}, \textquotedblleft{}country\textquotedblright{} (\emph{kwok}) (KANJIDIC2; Middle Chinese in Baxter's transcription).
\end{cousinweb}

\begin{grammarlens}[title={What you've built}]
One more word for getting around.
\end{grammarlens}

\subsection*{Your turn}
\begin{itemize}
\raggedright
  \item \emph{Say it:} \emph{nyūkoku}
  \item \emph{Say it:} \emph{nyūkoku}, once more
  \item \emph{Say it:} \emph{nyūkoku no toki, ryoken ga hitsuyō desu} — You need your passport when you enter the country.
  \item \emph{Recall:} say the Japanese for a passport, then the Japanese for customs, then say \emph{nyūkoku} again
\end{itemize}

\begin{tcolorbox}[breakable,colback=teal!4,colframe=teal!35!black,title={Before you move on}]
What does \ja{にゅうこく} mean? (\textquotedblleft{}Entry into a country\textquotedblright{}.) Say it once more.
\end{tcolorbox}
\section[\texorpdfstring{\ja{しゅっこく} (shukkoku)}{しゅっこく (shukkoku)}]{\ja{しゅっこく} (shukkoku) — departure from a country}

\label{lesson:JA-C386-shukkoku}



\begin{quote}
Before the new one: say the Japanese for customs, then the Japanese for entry into a country.
\end{quote}

\subsection*{You'll want to know: \ja{しゅっこく}}
\textbf{\ja{しゅっこく}} — \emph{shukkoku} — \textquotedblleft{}departure from a country\textquotedblright{}.

\begin{quote}
\textbf{\ja{しゅっこくの} \ja{まえに} \ja{おみやげを} \ja{かいました。}} — \emph{shukkoku no mae ni omiyage o kaimashita} — \textquotedblleft{}I bought souvenirs before leaving the country.\textquotedblright{}
\end{quote}

\begin{cousinweb}
\textbf{\ja{しゅっこく}} (\emph{shukkoku}) is Sino-Japanese (kango): two Chinese characters read with their on'yomi, readings borrowed from Chinese. They are \emph{shutsu}, \textquotedblleft{}exit\textquotedblright{} (Middle Chinese \emph{tsyhwit}) and \emph{koku}, \textquotedblleft{}country\textquotedblright{} (\emph{kwok}), joined as \emph{shukkoku} (KANJIDIC2; Middle Chinese in Baxter's transcription).
\end{cousinweb}

\begin{grammarlens}[title={What you've built}]
One more word for getting around.
\end{grammarlens}

\subsection*{Your turn}
\begin{itemize}
\raggedright
  \item \emph{Say it:} \emph{shukkoku}
  \item \emph{Say it:} \emph{shukkoku}, once more
  \item \emph{Say it:} \emph{shukkoku no mae ni omiyage o kaimashita} — I bought souvenirs before leaving the country.
  \item \emph{Recall:} say the Japanese for customs, then the Japanese for entry into a country, then say \emph{shukkoku} again
\end{itemize}

\begin{tcolorbox}[breakable,colback=teal!4,colframe=teal!35!black,title={Before you move on}]
What does \ja{しゅっこく} mean? (\textquotedblleft{}Departure from a country\textquotedblright{}.) Say it once more.
\end{tcolorbox}
\section[\texorpdfstring{\ja{とりだす} (toridasu)}{とりだす (toridasu)}]{\ja{とりだす} (toridasu) — to take something out; to produce}

\label{lesson:JA-C386-toridasu}



\begin{quote}
Before the new one: say the Japanese for entry into a country, then the Japanese for departure from a country.
\end{quote}

\subsection*{You'll want to know: \ja{とりだす}}
\textbf{\ja{とりだす}} — \emph{toridasu} — \textquotedblleft{}to take something out; to produce\textquotedblright{}.

\begin{quote}
\textbf{\ja{かばんから} \ja{りょけんを} \ja{とりだしました。}} — \emph{kaban kara ryoken o toridashimashita} — \textquotedblleft{}I took my passport out of my bag.\textquotedblright{}
\end{quote}

Polite \textbf{\ja{とりだします}}, \emph{toridashimasu}; a godan verb. Past (polite): \textbf{\ja{ちちが} \ja{かばんから} \ja{りょけんを} \ja{とりだしました。}}, \emph{chichi ga kaban kara ryoken o toridashimashita}.

\begin{cousinweb}
\textbf{\ja{とりだす}} (\emph{toridasu}) is a native Japanese word (wago), inherited rather than borrowed from Chinese, built from \emph{tori}, \textquotedblleft{}taking\textquotedblright{} (from \emph{toru}) and \emph{dasu}, \textquotedblleft{}to put out\textquotedblright{}; its kanji were attached later as kun'yomi (JMdict).
\end{cousinweb}

\begin{grammarlens}[title={What you've built}]
One more word for getting around.
\end{grammarlens}

\subsection*{Your turn}
\begin{itemize}
\raggedright
  \item \emph{Say it:} \emph{toridasu}
  \item \emph{Say it:} \emph{toridasu}, once more
  \item \emph{Say it:} \emph{kaban kara ryoken o toridashimashita} — I took my passport out of my bag.
  \item \emph{Recall:} say the Japanese for entry into a country, then the Japanese for departure from a country, then say \emph{toridasu} again
\end{itemize}

\begin{tcolorbox}[breakable,colback=teal!4,colframe=teal!35!black,title={Before you move on}]
What does \ja{とりだす} mean? (\textquotedblleft{}To take something out; to produce\textquotedblright{}.) Say it once more.
\end{tcolorbox}
